%! Author = Sujal Singh
%! Date = 10/23/23

% Preamble
\documentclass[11pt,a4paper,oneside]{scrarticle}

% Packages
\usepackage{mathptmx}
\usepackage{tabularray}
\usepackage{graphicx}
\usepackage{sectsty}
\usepackage{titlesec}
\usepackage{setspace}
\usepackage{romannum}
\usepackage[
    pdftitle={The Indian Government's Stance on Fundamental Rights in the Digital World},
    pdfsubject={The Indian Government's Stance on Fundamental Rights in the Digital World},
    pdfauthor={Sujal Singh},
    pdfdisplaydoctitle,
    hidelinks,
]{hyperref}
\usepackage[shortlabels]{enumitem}

\renewcommand*{\maketitle}{
    \pagenumbering{gobble}
    \begin{center}
        \vspace*{8em}
        \includegraphics[width=200pt]{../../tools/pdf/template/logo} \\[40pt]
        \textbf{\Huge Law Assignment} \\[40pt]
        \begin{table}[htb]
            \label{tab:name-slip}
            \begin{tblr}{rows={30pt},colspec={|X|X[3]|},hlines}
                \textbf{Course}            & Industrial Internet of Things \\
                \textbf{Batch}             & 1                             \\
                \textbf{Name}              & Sujal Singh                   \\
                \textbf{Enrollment Number} & XXXXXXXXXXXX                  \\
                \textbf{Roll Number}       & XX                            \\
            \end{tblr}
        \end{table}
    \end{center}
    \newpage
    \pagenumbering{arabic}
}

\linespread{1.5}
\titleformat{\section}{\Large\bfseries\centering}{\thesection.}{0.5em}{}

\newcommand{\link}[1]{\href{#1}{\url{#1}}}

% Document
\begin{document}
    \maketitle

    \begin{center}
        \textbf{\LARGE The Indian Government's Stance on Fundamental Rights in the Digital World.}
        \\[10pt]
    \end{center}
    \begin{flushright}
        \textit{--- Sujal Singh} \\[20pt]
    \end{flushright}

    %------------------------------------------------------------------------------------------------------------------%
    %------------------------------------------------------------------------------------------------------------------%

    \section*{Abstract}\label{sec:abstract}
    The Internet is one of mankind's greatest achievements and an open Internet, free of censorship and control by any
    authority, is crucial to avoid a dystopian future, but ever since it's inception, government authorities have
    employed various tactics to censor the information published via the World Wide Web.\ This report details the
    state of fundamental rights as they exist in the digital landscape of India and how the Indian government has often
    violated these rights to censor the Internet and use it as a tool for mass-surveillance.\ Such acts go against
    constitutional values as envisioned by the authors of the Indian Constitution, and people in relevant fields of
    studies should feel responsible to oppose the continued decline of the open Internet, and they should push for
    technical solutions that uphold constitutional values.


    %------------------------------------------------------------------------------------------------------------------%
    %------------------------------------------------------------------------------------------------------------------%


    \section{Introduction}\label{sec:introduction}

    The Indian Constitution provides six fundamental rights to all citizens of India, one of which is the right to
    freedom.\ According to article 19(1)(a):
    \begin{quote}
        \singlespacing
        (1)
        All citizens shall have the right---\\
        \vspace*{-15pt}
        \begin{enumerate}
        (a)
            to freedom of speech and expression;\\
            \ldots
        \end{enumerate}
    \end{quote}
    citizens of India are provided a fundamental right to express their views relating to any matter.\ Due to this
    fundamental right, citizens are allowed to offer criticism to the state, in fact, without criticism the government
    loses accountability, and yet, various articles offering such criticism published on the web often has the
    possibility of landing the author in jail.
    \\
    According to the Supreme Court of India's ruling in the
    \textit{Justice K. S. Puttaswamy (Retd.) \& Another. vs. Union Of India \& Others} case, the right to privacy is
    protected as an intrinsic part of the right to life and personal liberty under Article 21 and as a part of the
    freedoms guaranteed by Part \Romannum{3} of the Constitution.
    \\
    These fundamental rights are guaranteed for every citizen of India, but in reality, the Indian government itself has
    a history of openly violating these rights.\ The following sections describe how the Indian government tackles
    various issues relating to maintaining fundament rights on the Internet.


    %------------------------------------------------------------------------------------------------------------------%
    %------------------------------------------------------------------------------------------------------------------%


    \section{Censorship}\label{sec:censorship}

    Internet Service Providers (ISPs) are often ordered by the Department of Telecommunications (DoT) to block a large
    number of websites without any prior warnings or even without stating any reasons as to why a particular site was
    added to the blocklist.\ This blocklist may include websites which distribute any content unfavourable to the
    government, this violates the fundament right of speech and expression.\ The lack of transparency by the government
    on what websites are added to the blocklist is a violation of right to information and goes against the spirit of
    democracy, and often leads to legitimate websites being blocked.\ I have personally experienced two such instances;
    the content delivery network of a major Microsoft owned source-code hosting website, GitHub, was blocked for months
    by major ISPs such as Airtel, Jio, ACT Fibernet and many more.\ The website for the free and open source software
    VLC was also blocked under orders from the Ministry of Electronics and Information Technology.\ Moreover, the
    technical implementation of this blocklist inherently implies the violation of right to privacy as one can not block
    a website without first knowing what websites a person is visiting in the first place.\ Till now, ISPs chose to
    implement these orders by returning a bogus Internet address when a user requested the location of a blocked
    website from their servers, this implementation could be easily bypassed by simply directing encrypted queries to a
    server outside the jurisdiction of the Indian government, but now, large ISPs like Airtel and Jio have found more
    successful methods to implement the blocklist, they exploit an inherent flaw in the encryption protocol used on the
    web (TLS), wherein, the ISP can not know what contents on a website a user is accessing, but it can figure out what
    website they're accessing as that information is sent out in plain-text during the initial handshake between the
    server and the client.


    %------------------------------------------------------------------------------------------------------------------%
    %------------------------------------------------------------------------------------------------------------------%


    \section{Aadhar Card}\label{sec:aadhar-card}
    Ever since the launch of the Aadhar Card by the Unique Identification Authority of India (UIDAI), it has attracted a
    lot of criticism which has resulted in numerous Supreme Court rulings regarding Aadhar.\ The crux of the issue with
    Aadhar is the right to privacy, or the lack thereof, and the implications of a central database uniquely identifying
    every citizen.\ While the Supreme court has ruled that Aadhar cards shall not be made mandatory for any essential
    services, in reality, most government tasks have made the alternate options highly inconvenient virtually rendering
    Aadhar as a mandatory document.\ It's not unlikely that a large portion of the population, including government
    officials, might not even be aware that Aadhar is not a mandatory document.
    \\
    Most computer security experts have expressed concerns that a central database like Aadhar is a major
    security risk to its citizens, for it's not a question of if the database will suffer a breach but a question of
    when.\ And indeed, these predictions have come true numerous times, exposing highly sensitive information of Indian
    citizens to foreign attackers.\ The fact is that maintaining security on such a critical project would prove to be
    particularly hard even for big tech companies, let alone the largely technically incompetent Indian government.
    \\
    Even disregarding whether the government can maintain security on these systems, Aadhar serves as the link between
    disjoint datasets allowing the government to profile its citizens.\ This threat is particularly more prominent now
    in the age of artificial intelligence.


    %------------------------------------------------------------------------------------------------------------------%
    %------------------------------------------------------------------------------------------------------------------%


    \section{Virtual Private Network Providers}\label{sec:virtual-private-network-providers}

    The Indian Computer Emergency Response Team (CERT-In) recently started requiring all companies providing VPN
    services to keep log of the traffic flowing through their servers located in India.\ This resulted in several VPN
    companies removing their servers from India, as complying with such laws would go against the business model these
    companies operate on.\ This is evidence of the deteriorating state of the open Internet in India and had companies
    obliged, it would've been a violation of the right to privacy.\ But this law is not limited to VPN providers, it
    also states that any Virtual Private Server/Cloud service providers shall maintain an accurate log of the following
    information for a period of 5 years or longer even after a cancellation/withdrawal by the user:

    \begin{enumerate}[(a)]
        \singlespacing
        \item Validated names of subscribers/customers hiring the services
        \item Period of hire including dates
        \item IPs allotted to / being used by the members
        \item Email address and IP address and time stamp used at the time of registration / on-boarding
        \item Purpose for hiring services
        \item Validated address and contact numbers
        \item Ownership pattern of the subscribers / customers hiring services
    \end{enumerate}


    %------------------------------------------------------------------------------------------------------------------%
    %------------------------------------------------------------------------------------------------------------------%


    \section{Net Neutrality}\label{sec:net-neutrality}

    The Indian government has had a relatively good stance on the topic of net neutrality, while there is no legislation
    regarding this matter in India, the Department of Telecommunications (DoT) upon recommendations from the Telecom
    Regulatory Authority of India (TRAI) has barred ISPs from performing any from of data discrimination.\ Non-profit
    organizations like the Internet Freedom Foundation continue to fight regarding such issues against various concerned
    regulatory bodies of India.


    %------------------------------------------------------------------------------------------------------------------%
    %------------------------------------------------------------------------------------------------------------------%


    \section*{Conclusion}\label{sec:conclusion}
    India ranks 51\textsuperscript{st} in the Internet Freedom report, marked as \textit{``partly free''}, and I believe
    India is in fact in a relatively good position compared to most of our neighbouring countries, but we the citizens
    shall continue to hold the government accountable to counteract the worsening state of fundamental rights in the
    digital landscape of India.


    %------------------------------------------------------------------------------------------------------------------%
    %------------------------------------------------------------------------------------------------------------------%


    \newpage
    \section*{References}\label{sec:references}

    \begin{enumerate}
        \singlespacing
        \item \link{https://legislative.gov.in/constitution-of-india/}
        \item \link{https://wikipedia.org/wiki/Fundamental_rights_in_India}
        \item \link{https://wikipedia.org/wiki/Right_to_Privacy_verdict}
        \item \link{https://wikipedia.org/wiki/Information_Technology_Act,_2000}
        \item \link{https://wikipedia.org/wiki/Internet_censorship_in_India}
        \item \link{https://wikipedia.org/wiki/Aadhaar}
        \item \link{https://thewire.in/government/privacy-aadhaar-supreme-court}
        \item \link{https://www.cert-in.org.in/PDF/CERT-In_Directions_70B_28.04.2022.pdf}
        \item \link{https://wikipedia.org/wiki/Net_neutrality_in_India}
        \item \link{https://internetfreedom.in/}
        \item \link{https://freedomhouse.org/countries/freedom-net/scores}
    \end{enumerate}
\end{document}